% BHU PYQ -- Manually transcribed & error-corrected
% Paper: MTB-102 : Calculus
% Exam: B.Sc. (Hons.) Semester I Examination, 2016-17 (Old Course)

\documentclass[11pt,a4paper]{article}
\usepackage[a4paper,top=2.5cm,bottom=2.5cm,left=2.8cm,right=2.8cm,headheight=14pt]{geometry}
\usepackage{amsmath,amssymb}
\usepackage[T1]{fontenc}
\usepackage[utf8]{inputenc}
\usepackage{lmodern}
\usepackage{microtype}
\usepackage{fancyhdr}
\usepackage{enumitem}
\usepackage{hyperref}
\usepackage{lastpage}
\usepackage{parskip}

\hypersetup{colorlinks=true,urlcolor=black,linkcolor=black,
  pdftitle={MTB-102: Calculus -- BHU B.Sc. Sem I 2016-17 (Old Course)}}
\pagestyle{fancy}\fancyhf{}
\renewcommand{\headrulewidth}{0.4pt}\renewcommand{\footrulewidth}{0.4pt}
\fancyhead[L]{\small\textit{Banaras Hindu University}}
\fancyhead[R]{\small\textit{MTB-102: Calculus (Old Course)}}
\fancyfoot[C]{\small Page \thepage\ of \pageref{LastPage}}

\newlist{parts}{enumerate}{2}
\setlist[parts,1]{label=(\alph*),leftmargin=*,itemsep=5pt,topsep=3pt}
\setlist[parts,2]{label=(\roman*),leftmargin=*,itemsep=3pt,topsep=2pt}
\newcommand{\pts}[1]{\hfill\textit{[#1]}}

\begin{document}
\begin{center}
  {\large\bfseries BANARAS HINDU UNIVERSITY}\\[2pt]
  {\normalsize B.Sc. (Hons.) --- Semester I Examination, 2016-17 (Old Course)}\\[6pt]
  \rule{0.8\linewidth}{0.5pt}\\[6pt]
  {\large\bfseries Paper No. MTB-102: Calculus}\\[4pt]
  {\normalsize Time: 3 Hours \hfill Full Marks: 70}
\end{center}
\rule{\linewidth}{0.4pt}
\smallskip
\noindent\textit{Instructions: Answer \textbf{five} questions including Question~1 which is compulsory. Figures in right-hand margin indicate marks.}
\bigskip

\noindent\textbf{Question 1.} \textit{(Compulsory --- 2 marks each.)}
\begin{parts}
  \item Give an example of a sequence which is bounded but not convergent.
  \item By using $\epsilon - \delta$ definition, prove that $\lim_{x\to 2} (x^2 + 2) = 6$.
  \item Show that the function defined by
  \[
  f(x) = \begin{cases} \frac{|x-a|}{x-a}, & x \neq a \\ 1, & x = a \end{cases}
  \]
  is not continuous at $x = a$. Write the type of discontinuity at $x = a$.
  \item If $2^x + 2^y = 2^{x+y}$, then find the value of $\frac{dy}{dx}$.
  \item If $y = \sin^2 x$, then find $y_n$.
  \item If $\cos x = a_0 + a_1 x + a_2 x^2 + a_3 x^3 + \dots$, then find the value of $a_3$.
  \item Expand $f(x) = 2x^3 + 5x^2 + x - 1$ in powers of $(x - 1)$.
  \item Verify Rolle's theorem for the function $f(x) = x^2 - 4x + 3$ on $[1,3]$.
  \item Find all asymptotes of the curve $xy - y = 0$.
  \item Find the nature of the multiple point $(0, 0)$ on the curve $x^3 + y^3 - 3axy = 0$.
  \item Find the value of $\lim_{n\to\infty} \left[ \frac{1}{n} + \frac{1}{n+1} + \frac{1}{n+2} + \dots + \frac{1}{2n} \right]$.
\end{parts}

\medskip\rule{\linewidth}{0.2pt}

\medskip\noindent\textbf{Question 2.}
\begin{parts}
  \item Let $\dots \lim_{x\to a} f(x) = l$ and $\lim_{x\to a} g(x) = m$. Prove that $\lim_{x\to a} f(x)g(x) = lm$. \pts{7}
  \item Examine the continuity of the function $f(x) = [x]$, where $[x]$ is the greatest integer function, at $x = 0.5$ and $x = 1$. \pts{5}
\end{parts}

\medskip\rule{\linewidth}{0.2pt}

\medskip\noindent\textbf{Question 3.}
\begin{parts}
  \item State and prove Lagrange's mean value theorem. Also write its geometrical meaning. \pts{6}
  \item Prove that for all finite values of $x$:
  \[
  e^x \sin x = x + x^2 + \frac{2x^3}{3!} - \frac{4x^5}{5!} + \dots
  \] \pts{6}
\end{parts}

\medskip\rule{\linewidth}{0.2pt}

\medskip\noindent\textbf{Question 4.}
\begin{parts}
  \item Show that every differentiable function is continuous, but the converse is not true. \pts{5}
  \item Test the differentiability of the function
  \[
  f(x) = \begin{cases} 2+x, & x < 1 \\ 4-x, & x \ge 1 \end{cases}
  \]
  at $x = 1$. \pts{4}
  \item By using Cauchy's mean value theorem, show that:
  \[
  \frac{\sin\beta - \sin\alpha}{\cos\alpha - \cos\beta} = \cot\theta, \quad 0 < \alpha < \theta < \beta < \frac{\pi}{2}
  \] \pts{3}
\end{parts}

\medskip\rule{\linewidth}{0.2pt}

\medskip\noindent\textbf{Question 5.}
\begin{parts}
  \item State and prove Leibniz's theorem. \pts{5}
  \item If $y = \cos(m \sin^{-1} x)$, $m \in \mathbb{N}$, then show that:
  \[
  (1-x^2) y_{n+2} - (2n+1) x y_{n+1} + (m^2 - n^2) y_n = 0
  \]
  and hence find $y_n(0)$. \pts{7}
\end{parts}

\medskip\rule{\linewidth}{0.2pt}

\medskip\noindent\textbf{Question 6.}
\begin{parts}
  \item Find all the asymptotes of the curve:
  \[
  x^3 + 2x^2y - xy^2 - 2y^3 + xy - y^2 + 1 = 0
  \] \pts{5}
  \item Trace the curve $y(4-x^2) = 4x$. \pts{7}
\end{parts}

\medskip\rule{\linewidth}{0.2pt}

\medskip\noindent\textbf{Question 7.}
\begin{parts}
  \item Trace the curve $r = a(1 + \cos\theta)$, $a > 0$. \pts{6}
  \item By using the definition of definite integral as the limit of a sum, evaluate $\int_a^b \cos x \, dx$. \pts{6}
\end{parts}

\vfill
\rule{\linewidth}{0.4pt}
\begin{center}\small
  \href{https://bhu-exam.vercel.app/}{Click here to visit our website}\\[4pt]
  \href{https://me-aryan.vercel.app/}{Click here to visit our contributor}\\[4pt]
  \href{https://quashan-me.vercel.app/}{Click here to visit our contributor}
\end{center}

\end{document}
